\documentclass[10pt,a4paper]{amsart}
\usepackage[margin=19mm]{geometry}
\usepackage{amsmath,amssymb,mathtools}
\usepackage[scr=rsfs,cal=euler]{mathalfa}
\usepackage{xfrac}
\usepackage[unicode,hidelinks,bookmarks=false]{hyperref}
\newcommand{\ie}{i.e.\ }
\newcommand{\cf}{cf.\ }
\newcommand{\resp}{respectively\ }
\newcommand{\Subsection}{\subsection}
\providecommand{\nobreakdash}{\nobreak\hbox{-}}
\setlength{\emergencystretch}{2em}
\multlinegap0pt
\usepackage{xcolor}
\usepackage{amssymb}
\usepackage[normalem]{ulem}
\usepackage[shortlabels]{enumitem}
\usepackage{caption}
\newtheorem{cor}{Corollary}
\usepackage{cleveref}
\crefname{cor}{Corollary}{Corollarys}
\newcommand{\IC}{\mathbb C}
\title{Equivariant Chevalley, Giambelli, and Monk Formulae for the Peterson Variety}
\author{Rebecca Goldin, Rahul Singh}
\date{Source: arXiv:2111.15663v2 (CC BY 4.0)}
\begin{document}
\maketitle

\noindent\emph{Source and licence.} Equivariant Chevalley, Giambelli, and Monk Formulae for the Peterson Variety. Version arXiv:2111.15663v2 (CC BY 4.0), distributed by arXiv under the Creative Commons Attribution 4.0 International licence, http://creativecommons.org/licenses/by/4.0/, as recorded in the arXiv licence metadata for this exact version and shown on the abstract page https://arxiv.org/abs/2111.15663v2. Full licence notice: \texttt{licenses/arxiv-2111.15663-CC-BY-4.0.txt}.\par\smallskip

\noindent\textit{Editorial scope.} Counted unit: the article's cor cor:twistedSection (source0.tex lines 734--753). The article's complete original proof of it is delivered with it (source0.tex lines 754--790, one \texttt{proof} environment of 37 lines). No mathematics is written, repaired or re-derived: the statement and the proof are the article's own lines, in the article's own notation, and no retained line is substituted or re-flowed.  Retained uncounted as the source-local context the proof needs: lines 720--725.  Nothing was omitted from any retained span, so the package carries no omission record.  No retained line contains a citation, so the wrapper carries no bibliography and no bibliography entry.  No retained line contains a figure environment, an image-inclusion command or drawing code, so no image is delivered and the package declares no asset dependency.  Editorial formatting only, in addition: an AMS \texttt{amsart} preamble that restates the article's own declarations and macros used by the retained lines, namely the environments \texttt{cor} on separate counters, together with the standard packages those lines need.  The retained environments print in this document as Corollary 1 in order of appearance, whereas the article prints the same items with the numbering of its own full document.  The complete native source of the article as distributed in the arXiv e-print for this exact version is bundled unchanged as \texttt{sources/119090/source0.tex}.
\begin{cor}
\label{cor:EulerClass}
If $\mathcal V$ is a $T$-equivariant vector bundle on a $T$-scheme $X$,
and $s$ is a $T$-invariant regular section of $\mathcal V$ with zero-scheme $Y$,
then $[Y]_T=e^T(\mathcal V)\cap[X]_T$.
\end{cor}

\begin{cor}
\label{cor:twistedSection} 
Let $\mathcal V\to X$ be  a $T$-equivariant vector bundle. 
For a character $\lambda$ of $T$, 
let $s$ be a regular section of $\mathcal V$ that lies in the $\lambda$-weight space, i.e.,
$(z\cdot s)=\lambda(z) s $ for all $z\in T$.
%\begin{align*}
%&&(z\cdot s)=\lambda(z) s && \forall z\in T.
%\end{align*}
The zero scheme $Z(s)$ of $s$ is $T$-invariant, and we have
%Let $Y$ be the zero scheme of $s$, and note that $[Y]$ is $T$-invariant.
\begin{equation*}
[Z(s)]_T=e^{T}\left(\mathcal V\otimes\underline\IC_{-\lambda}\right)\cap[X]_T.
\end{equation*}
If $\mathcal V$ admits a filtration with $T$-equivariant line bundle quotients $\{\mathcal L_i\}$,
we have
\begin{equation*}
[Z(s)]_T=\prod\left(c_{1}^{T}(\mathcal L_i)-\lambda\right)\cap[X]_T.
\end{equation*}
\end{cor}

\begin{proof}
%Let $\underline{\IC}_{-\lambda}=X\times\IC\to X$ be the (geometrically trivial) equivariant line bundle
%with $T$-action given by $z(x,v)=(zx,\lambda(z^{-1})v)$.
If $\lambda$ is non-trivial, the section $s$ is not invariant.
Observe however that the section $r=s\otimes 1$
of the vector bundle $\mathcal V\otimes\underline\IC_{-\lambda}$ is $T$-invariant, since
%$t s(t^{-1} x)=\lambda(t) s(x)\implies s(t^{-1}x)=\lambda(t) t^{-1} s(x)$.
%\begin{equation*}
%r(zx)=s(zx)\otimes 1=\lambda(z^{-1})zs(x)\otimes 1=zs(x)\otimes \lambda(z^{-1})=z(s(x)\otimes 1)=zr(x),
%\end{equation*}
\begin{align*}
r(z x)&=s(z x)\otimes 1=\left(zz^{-1}s(zx)\right)\otimes 1\\
&= \left(z\left( \left(z^{-1}\cdot s\right)(x) \right)\right) \otimes 1 \\
&= \left(z\left( \lambda(z^{-1})s(x) \right)\right)\otimes 1\\
& = \left(\lambda(z^{-1})zs(x)\right)\otimes 1 \\
&=zs(x)\otimes \lambda(z^{-1})=z(s(x)\otimes 1)=zr(x).
\end{align*}
%But what we want is not this. Instead, we show that the section is {\em invariant}, i.e.
%\begin{align*}
%(z\cdot r)(x) &= z\cdot (s\otimes 1)(x) = (z\cdot s)(x)\otimes (z\cdot 1)\\
%& = (\lambda(z)s(x) \otimes \lambda(z^{-1}) = s(x)\otimes 1 = r(x).
%\end{align*}
Further, $r$ has the same zero scheme as $s$, i.e., $Z(r)=Z(s)$.
%Then $s$ induces an equivariant section of the twisted bundle $\mathcal V\otimes V_\chi$.
Hence the first equality follows from \cref{cor:EulerClass}.

Suppose $\mathcal V$ admits a filtration by line bundles $\{\mathcal L_i\}$.
Then $\mathcal V\otimes\underline\IC_{-\lambda}$
admits a filtration by line bundles $\{\mathcal L_i\otimes\underline\IC_{-\lambda}\}$.
Applying the Whitney splitting principle, we have
\begin{align*}
e^T\left(\mathcal V\otimes\underline\IC_{-\lambda}\right)
=\prod\left(c_1^{T}\left(\mathcal L_i\otimes\underline\IC_{-\lambda}\right)\right)
=\prod\left(c_1^{T}(\mathcal L_i)-\lambda\right),
\end{align*}
from which the second equality follows.
\end{proof}

\end{document}
